\appendix

\section{Netlists dos circuitos}
\label{ap:netlist}

Netlists exatamente como foram simuladas (LTspice 26 em Docker,
\texttt{./claw-spice sim run}). O arquivo \texttt{vos\_comp.inc} é
gerado automaticamente com o $V_{OS}$ medido no item 3 e define o
parâmetro \texttt{VOSC} usado nos itens 4 a 6.

\subsection*{Item 1: \texttt{item1\_bias.cir}}
{\small
\begin{verbatim}
* RS1/S1 item 1 - Fig. 1(a): corrente de polarizacao de entrada IB+ (RI=1Meg, RF=0)
* OP07 da biblioteca embarcada do LTspice (LTC.lib). Alimentacao simetrica +/-15 V
* explicita (o Micro-Cap a inclui automaticamente; no LTspice ela deve ser declarada).
* Com RF=0 o amplificador e um seguidor: V(out) = V(inp) = -IB+ * RI.
* IB+ = -I(RI) (corrente que entra no pino nao inversor vinda do terra atraves de RI).
* O LTspice 26 nao aceita resistencia zero em netlist; RF=1uOhm implementa o
* curto do roteiro com erro desprezivel (IB- * 1uOhm ~ 1e-15 V).
.lib LTC.lib
VCC vcc 0 15
VEE vee 0 -15
RI inp 0 1Meg
RF out inn 1u
XU1 inp inn vcc vee out OP07
.op
.end
\end{verbatim}
}

\subsection*{Item 2: \texttt{item2\_ios.cir}}
{\small
\begin{verbatim}
* RS1/S1 item 2 - Fig. 1(a) com RF=1Meg: correntes IB- e de offset Ios
* Mesma topologia do item 1, agora com RF=1Meg na realimentacao.
* IB- = I(RF) (corrente que entra no pino inversor vinda da saida atraves de RF).
* V(out) = RF*IB- - RI*IB+  =>  com RI=RF=1Meg, V(out) ~ -Ios*1Meg.
* Ios = |IB+ - IB-| (corrente de offset de entrada).
.lib LTC.lib
VCC vcc 0 15
VEE vee 0 -15
RI inp 0 1Meg
RF out inn 1Meg
XU1 inp inn vcc vee out OP07
.op
.end
\end{verbatim}
}

\subsection*{Item 3: \texttt{item3\_vos.cir}}
{\small
\begin{verbatim}
* RS1/S1 item 3 - Fig. 1(b): tensao de offset de entrada VOS (malha aberta)
* Vin DC no pino nao inversor, pino inversor aterrado, sem realimentacao.
* Varredura DC de -5 mV a +5 mV com passo de 1 uV (10001 pontos).
* VOS = valor de Vin para o qual V(out) cruza 0 V na curva V(out) x Vin.
.lib LTC.lib
VCC vcc 0 15
VEE vee 0 -15
VIN vin 0 0
XU1 vin 0 vcc vee out OP07
.dc VIN -5m 5m 1u
.meas DC vos FIND V(vin) WHEN V(out)=0 CROSS=1
.end
\end{verbatim}
}

\subsection*{Item 4: \texttt{item4\_avd.cir}}
{\small
\begin{verbatim}
* RS1/S1 item 4 - Fig. 1(c): ganho DC diferencial Avd
* Fonte VOSC em serie com o pino nao inversor cancela o offset medido no item 3,
* mantendo o amplificador na regiao linear em torno de Vin=0.
* Avd = max |d V(out) / d Vin| (derivada numerica calculada em analyze.py,
* equivalente ao comando dd(u) do Micro-Cap). As .meas abaixo dao uma
* estimativa pela inclinacao entre -10 uV e +10 uV.
.lib LTC.lib
.include vos_comp.inc
VCC vcc 0 15
VEE vee 0 -15
VIN vin 0 0
VOS inp vin {VOSC}
XU1 inp 0 vcc vee out OP07
.dc VIN -5m 5m 1u
.meas DC v_m10u FIND V(out) AT=-10u
.meas DC v_p10u FIND V(out) AT=10u
.meas DC avd_sec PARAM (v_p10u-v_m10u)/20u
.end
\end{verbatim}
}

\subsection*{Item 4, segundo caso: \texttt{item5\_acm.cir}}
{\small
\begin{verbatim}
* RS1/S1 item 5 - Fig. 1(d): ganho DC de modo comum Acm
* Os dois pinos de entrada sao excitados juntos por Vin (modo comum); a fonte
* VOSC em serie com o pino nao inversor cancela o offset para evitar saturacao.
* Varredura de -3 V a +3 V (passo 1 mV). Acm = max |d V(out) / d Vin|.
* CMRR = 20*log10(Avd/|Acm|) e calculada em analyze.py com o Avd do item 4.
.lib LTC.lib
.include vos_comp.inc
VCC vcc 0 15
VEE vee 0 -15
VIN vin 0 0
VOS inp vin {VOSC}
XU1 inp vin vcc vee out OP07
.dc VIN -3 3 1m
.meas DC v_m1 FIND V(out) AT=-1
.meas DC v_p1 FIND V(out) AT=1
.meas DC acm_sec PARAM (v_p1-v_m1)/2
.end
\end{verbatim}
}

\subsection*{Item 6: \texttt{item6\_slew.cir}}
{\small
\begin{verbatim}
* RS1/S1 item 6 - Fig. 1(e): taxa de inflexao (slew rate)
* Seguidor unitario (saida ligada ao pino inversor) com degrau abrupto de
* Vin: -15 V -> +15 V em t=20 us com trise=1 ps (aprox. 0 ps do roteiro).
* SR = max |d V(out)/dt| no flanco de subida (ddt numerico em analyze.py).
* As .meas abaixo dao a estimativa classica de datasheet entre -10 V e +10 V.
.lib LTC.lib
.include vos_comp.inc
VCC vcc 0 15
VEE vee 0 -15
VIN vin 0 PULSE(-15 15 20u 1p 1p 460u 1m 1)
VOS inp vin {VOSC}
XU1 inp out vcc vee out OP07
.tran 0 200u 0 100n
.meas TRAN t_lo WHEN V(out)=-10 RISE=1
.meas TRAN t_hi WHEN V(out)=10 RISE=1
.meas TRAN sr_avg PARAM (20/(t_hi-t_lo))/1e6
.end
\end{verbatim}
}

\subsection*{Teste extra (item 7): \texttt{item7\_slew\_ada.cir}}
{\small
\begin{verbatim}
* RS1 - teste extra: slew rate de um AmpOp moderno (ADA4610) para comparar
* com o OP07 do item 6. Mesmo seguidor unitario e mesmo degrau abrupto de
* Vin: -15 V -> +15 V em t=20 us com trise=1 ps. Sem compensacao de VOS
* (o VOS do ADA4610 e da ordem de microvolts e nao afeta a medida de slew).
* Modelo ADA4610 embarcado no LTspice (ADA4610.lib).
.lib ADA4610.lib
VCC vcc 0 15
VEE vee 0 -15
VIN vin 0 PULSE(-15 15 20u 1p 1p 460u 1m 1)
XU1 vin out vcc vee out ADA4610
.tran 0 60u 0 10n
.meas TRAN t_lo WHEN V(out)=-10 RISE=1
.meas TRAN t_hi WHEN V(out)=10 RISE=1
.meas TRAN sr_avg PARAM (20/(t_hi-t_lo))/1e6
.end
\end{verbatim}
}

\subsection*{Compensação de offset: \texttt{vos\_comp.inc}}
{\small
\begin{verbatim}
* GERADO por analyze.py --phase vos: compensacao com o VOS medido no item 3.
.param VOSC=-3.035132649e-05
\end{verbatim}
}

\section{Folha de dados do OP07 (páginas de especificações)}
\label{ap:datasheet}

Reprodução das páginas de especificações elétricas da folha de dados do
OP07 (Analog Devices, Rev. G), usadas como referência na
Tabela \ref{tab:medidas}: características do OP07E (página 3),
desempenho dinâmico do OP07E e características do OP07C (página 4) e
desempenho dinâmico do OP07C (página 5).

\begin{figure}[p]
  \centering
  \fbox{\includegraphics[page=3,width=0.9\textwidth]{datasheet-op07.pdf}}
  \caption{Folha de dados do OP07 (Rev. G), página 3: características
  elétricas do OP07E ($V_{OS}$, $I_{OS}$, $I_B$, CMRR, $A_{VO}$).}
  \label{fig:ds-p3}
\end{figure}

\begin{figure}[p]
  \centering
  \fbox{\includegraphics[page=4,width=0.9\textwidth]{datasheet-op07.pdf}}
  \caption{Folha de dados do OP07 (Rev. G), página 4: desempenho
  dinâmico do OP07E (slew rate) e características elétricas do OP07C.}
  \label{fig:ds-p4}
\end{figure}

\begin{figure}[p]
  \centering
  \fbox{\includegraphics[page=5,width=0.9\textwidth]{datasheet-op07.pdf}}
  \caption{Folha de dados do OP07 (Rev. G), página 5: características de
  saída e desempenho dinâmico do OP07C.}
  \label{fig:ds-p5}
\end{figure}

\section{Reprodução do experimento}
\label{ap:reproducao}

\begin{verbatim}
# imagem Docker com LTspice (uma vez)
./claw-spice build-prebuilt

# folha de dados (analog.com bloqueia clientes nao-navegador;
# o script usa o espelho do Internet Archive)
scripts/models/fetch-op07-datasheet.sh

# simulacoes + medidas + graficos
experiments/rs1-nao-idealidades-op07/run_all.sh

# relatorio
./claw-spice report rs1-nao-idealidades-op07
\end{verbatim}

O cálculo das derivadas, a extração das medidas e a geração dos
gráficos estão no script \texttt{analyze.py} do experimento. Os
resultados numéricos são gravados em \texttt{results.json}.
